\chapter{Bewegungsbeschreibung und Bewegungsanalysesysteme}
\label{chap:Bewegungsanalyse}
Im Folgenden Kapitel wird ein Überblick über die Bewegungsbeschreibung und Modellierung gegeben und optische Bewegungsanalysesysteme vorgestellt. Am Ende des Kapitels wird das zu entwickelnde System kurz skizziert.

\section{Bewegungsbeschreibung}
\label{sec:anatomy}
Die menschliche Anatomie kann als eine kinematische Kette beschrieben werden. Eine kinematische Kette sind kinematische Elemente, in diesem Fall Knochen, die über Gelenke miteinander verbunden sind. 

\subsection{Handanatomie}
\label{subsec:Gelenkzuordnung}
In der Anatomie gibt es verschiedene Gelenkarten zwischen den unterschieden werden. In dieser Arbeit soll ein System für die Analyse von Handbewegungen entwickelt werden. Es werden die Gelenke einer Hand vorgestellt. Dies sind Kugelgelenke, Elipsoidgelenke, Sattelgelenke und Scharniergelenke.

\begin{figure}[h]
    \centering
    \begin{overpic}[width=0.75\textwidth,trim = 0 0 0 0]{02_grundlagen/images/Gelenkarten.pdf}
        \put(15,50){1}
        \put(60,50){2}
        \put(15,20){3}
        \put(60,20){4}
    \end{overpic}
    \caption{Schematische Darstellung der Gelenke mit den Bewegungsrichtungen: 1 Elipsoidgelenk, 2 Sattelgelenk, 3 Kugelgelenk, 4 Scharniergelenk.(Abbildungen der Gelenke ist aus \cite{appell_funktionelle_2008} übernommen)}
    \label{fig:UebersichtGelenke}
\end{figure}

Kugelgelenke besitzen insgesamt drei Freiheitsgrade, Elipsoidgelenke und Sattelgelenke jeweils zwei Freiheitsgrade und Scharniergelenke einen Freiheitsgrad. Die möglichen Bewegungen sind in Abbildung \ref{fig:UebersichtGelenke} dargestellt.


Die Hand kann modellhaft als eine kinematische Kette bestehend aus 21 kinematischen Elemente, die mit unterschiedlichen Gelenken miteinander verbunden sind. Die Kette beginnt am Unterarm, der über ein Elipsoidgelenk (Handgelenk) mit der Hand verbunden ist. Dieses Gelenk hat zwei Freiheitsgrade. Vom Handgelenk gehen fünf Unterketten aus, eine für jeden Finger. Die Daumenkette führt zu einem Sattelgelenk und ermöglicht Bewegungen in zwei Richtungen. Die vier Fingerketten führen anatomisch gesehen zu Kugelgelenken. Diese können aufgrund der Bänder und Sehnen nur in zwei Richtungen bewegt werden, die Adduktion/Abbduktion und die Flexion/Extension. Die Rotation um die Achse des kinematischen Elements ist nur passiv möglich, also durch Krafteinwirkung von außen \cite{appell_funktionelle_2008}. Diese fünf Gelenke werden in der Medizin \textit{Metakarpophalangealgelenke} (MCP) genannt \cite{appell_funktionelle_2008}. 

In den fünf Ketten folgen dann jeweils zwei Scharniergelenke. Das erste Scharniergelenk in jeder Kette wird in der Medizin \textit{Proximale Interphalangealgelenk} (PIP) und das letzte Gelenk der Kette wird \textit{Distale Interphalangealgelenk} (DIP) genannt. Die Scharniergelenke können jeweils eine Flexion/Extension durchführen. Die kinematische Kette ist in Abbildung \ref{fig:Gelenkszuordnung} dargestellt.

\begin{figure}[h]
    \centering
    \begin{overpic}[width=0.5\textwidth]{02_grundlagen/images/Gelenkzuordnung.pdf}
        \put(38,62){MCP}
        \put(25,20){PIP}
        \put(8,42.75){DIP}
    \end{overpic}
    \caption{Zuordnung der schematischen Gelenke zu den realen Gelenken einer Hand. Sattelgelenke sind in blau, Scharniergelenk in gelb, Elipsoidgelenke in rot und Kugelgelenke in grün dargestellt.}
    \label{fig:Gelenkszuordnung}
\end{figure} 


\subsection{Modellierung}
\label{subsec:ModellingJoints}
Die Anatomie der menschlichen Hand wird in dieser Arbeit vereinfacht als kinematische Kette beschrieben. Kinematische Ketten bestehen aus kinematischen Elementen, die über unterschiedliche Gelenke miteinander verbunden sind. Die kinematischen Elemente sind am Beispiel der Hand die Knochen, in der Modellierung werden diese als Zylinder mit der entsprechenden Dicke angenommen. Die Gelenke werden als Sattelgelenke angenommen, da bei der Hand Gelenke mit maximal zwei Freiheitsgraden vorkommen. Die kinematischen Elemente werden mit der Variable $k$ nummeriert, die Gelenkswinkel gehören jeweils zum vorherigen Element. Das Modell ist in Abbildung\,\ref{fig:ModellierungBewegung} skizziert.

\begin{figure}[h]
    \centering
    \begin{overpic}[width=\textwidth]{02_grundlagen/images/Def_ModellDrehung.png}
        \put(56,65){$y$}
        \put(41.5,56.25){$z$}
        \put(80,68){$k$}
        \put(20,68){$k+1$}
    \end{overpic}
    \caption{Modellierung einer kinematischen Kette: Die Abbildung zeigt eine kinematische Kette bestehend aus zwei Elementen mit dem zugehörigen Koordinatensystem. Unten ist exemplarisch ein Finger mit der zugehörigen Modellierung dargestellt.}
    \label{fig:ModellierungBewegung}
\end{figure}

Jedes kinematische Element und Gelenk besitzt ein lokales Koordinatensystem. Dieses ist wie in Abbildung\,\ref{fig:ModellierungBewegung} skizziert. Die Symetrieachse des Elements $\vec{e}\tindex{kin,$k$}$ definiert die $z$-Achse. Die $x$- und $y$- Achsen werden durch die Drehachsen des folgenden Gelenks definiert. Die Drehung um die $y$-Achse wird mit dem Zenitwinkel $\theta\tindex{Beu,$k$}$ beschrieben und hat einen Wertebereich von $-\frac{\pi}{2} < \theta\tindex{Beu,$k$} < \frac{\pi}{2}$. Die Drehung um die $x$-Achse wird mit dem Polarwinkel $\phi\tindex{Sp,$k$}$ beschrieben und hat einen Wertebereich von $-\pi < \phi\tindex{Sp,$k$} < \pi$. 

Mit den Winkeln des Gelenks lässt sich das Koordinatensystem des folgenden Elements beschreiben. Die $z$-Achse des folgenden Elements $\vec{e}\tindex{$z,k+1$} = \vec{e}\tindex{kin,$k$}$ ergibt sich dann durch eine Drehung um $\vec{e}\tindex{$x,k$}$ und $\vec{e}\tindex{$y,k$}$. Die Drehung wird mit Quaternionen beschrieben, diese werden aus der Rotationsachse $\vec{e}\tindex{rot}$ und dem Winkel $\phi$ zusammengesetzt 

\begin{equation}
    \bvec{q} = \begin{pmatrix}
        \cos(\frac{\phi}{2}) \\
        \sin(\frac{\phi}{2})\,ex\tindex{rot} \\
        \sin(\frac{\phi}{2})\,ey\tindex{rot} \\
        \sin(\frac{\phi}{2})\,ez\tindex{rot} 
    \end{pmatrix}
\end{equation}
Die Beugung und die Spreizung werden mit den folgenden Quaternionen beschrieben
\begin{equation}
    \bvec{q}\tindex{Sp,k} = \begin{pmatrix}
        \cos(\frac{\phi\tindex{Sp,$k$}}{2}) \\
        \sin(\frac{\phi\tindex{Sp,$k$}}{2})\,ex\tindex{$x,k$} \\
        \sin(\frac{\phi\tindex{Sp,$k$}}{2})\,ey\tindex{$x,k$} \\
        \sin(\frac{\phi\tindex{Sp,$k$}}{2})\,ez\tindex{$x,k$} 
    \end{pmatrix},
    \bvec{q}\tindex{Beu,k} = \begin{pmatrix}
        \cos(\frac{\theta\tindex{Beu,$k$}}{2}) \\
        \sin(\frac{\theta\tindex{Beu,$k$}}{2})\,ex\tindex{$y,k$} \\
        \sin(\frac{\theta\tindex{Beu,$k$}}{2})\,ey\tindex{$y,k$} \\
        \sin(\frac{\theta\tindex{Beu,$k$}}{2})\,ez\tindex{$y,k$} 
    \end{pmatrix}
\end{equation}
und aufmultipliziert zur Gesamtbewegung des Gelenks $\bvec{q}\tindex{Gel}$
\begin{equation}
    \bvec{q}\tindex{Gel,$k$} = \bvec{q}\tindex{Sp,$k$} \circ \bvec{q}\tindex{Beu,$k$}. 
\end{equation} 
Die Achsen des folgenden Elements werden wie folgt bestimmt

\begin{equation}
    \vec{e}\tindex{$z,k+1$} = \bvec{q}\tindex{Gel,$k$}\,\cdot \vec{e}\tindex{$z,k$} \cdot \,\bvec{q}\tindex{Gel,k}^{-1}.
\end{equation}

In der Nulllage, $\theta\tindex{Beu}=0,\phi\tindex{Sp}=0$, ergibt sich, dass die Symmetrieachsen miteinander übereinstimmen und $\vec{e}\tindex{$z,k+1$} = \vec{e}\tindex{$z,k$}$. Die $x$- und $y$-Achsen jedes Elements $\vec{e}\tindex{$x,k,0$},\vec{e}\tindex{$y,k,0$}$ werden vor der Laufzeit bestimmt, die gesamte kinematische Kette ist dafür in der Nulllage. Aus $\bvec{q}\tindex{Gel,$k$}$ und $\vec{e}\tindex{$x,k+1,0$},\vec{e}\tindex{$y,k+1,0$}$, werden die aktuellen $x$- und $y$-Achsen bestimmt

\begin{equation}
    \vec{e}\tindex{$x,k+1$} = \bvec{q}\tindex{Gel,$k$}\,\cdot \vec{e}\tindex{$x,k+1,0$} \cdot \,\bvec{q}\tindex{Gel,$k$}^{-1},
\end{equation}

\begin{equation}
    \vec{e}\tindex{$y,k+1$} = \bvec{q}\tindex{Gel,$k$}\cdot\, \vec{e}\tindex{$y,k+1,0$} \cdot \,\bvec{q}\tindex{Gel,$k$}^{-1}.
\end{equation}

Bei kinematischen Ketten mit mehr als einem Gelenk müssen die Bewegungen aller Gelenke aufmultipliziert werden.


\section{Konventionelle Systeme}
\label{sec:ConventionalSystems}
In der Filmindustrie, im medizinischen oder vielen anderen Bereichen werden bereits Systeme zur Bewegungsanalyse verwendet. Diese basieren in den meisten Fällen auf optischer Sensorik. Für Handbewegungen im speziellen werden in modernen Systemen auch magnetische Sensoren verwendet. Dabei gibt es bereits kommerzielle Systeme, als auch Systeme die in einem sehr frühen Entwicklungsstadium sind. Die Funktionsweise der Systeme wird hier kurz vorgestellt.

\subsection{Optische Systeme}
\label{subsec:opticalSystems}

Optische Systeme für die Bewegungsanalyse gelten als der Goldstandard. Dabei gibt es unterschiedliche Ansätze, wir verwenden ein System, der optische Marker in Kombination mit Infrarotlicht verwendet. Optische Marker sind kugelförmige Objekte, deren Oberfläche so beschichtet ist, dass Infrarotlicht gut reflektiert wird. Die Kameras haben eine ringförmige Infrarotlichtquelle, mit der die Marker angeleuchtet werden. Diese beiden Bestandteile sind in Abbildung\,\ref{fig:MaterialOMC} gezeigt.

\begin{figure}[h]
    \centering
    \subfloat[Kamera\,\cite{noauthor_miqus_nodate}]{
        \begin{overpic}[width=0.35\textwidth,trim = 0 0 0 0]{02_grundlagen/images/Optical_camera.jpg}
        \end{overpic}
        \label{fig:Camera}
    }
    \subfloat[Optischer Marker\,\cite{noauthor_spherical_nodate}]{
        \begin{overpic}[width=0.35\textwidth,trim = 0 0 0 0]{02_grundlagen/images/optical_marker.jpg}
        \end{overpic}
        \label{fig:Opticalmarker}
    }
    \caption{Optisches System: (a) zeigt beispielhaft eine Kamera für die Bewegungsanalyse. (b) zeigt einen optischen Marker, dessen Bewegung von einem Kamerasystem aufgenommen werden kann.}
    \label{fig:MaterialOMC}
\end{figure}

Eine beliebige Anzahl Kameras wird im Raum so aufgestellt, dass der interessierte Bereich möglichst von allen Kameras eingesehen werden kann. Dies ist in Abbildung\,\ref{fig:CameraPositioning} dargestellt.

In einer Kalibrierungsroutine werden mithilfe eines Referenzobjektes, bestehend aus mehreren optischen Markern, die Position und Orientierung der Kameras bestimmt\,\cite{qualisys_ab_qtm-user-manual_nodate}.

Während der Laufzeit nehmen die Kameras synchronisiert Bilder auf und die Auswertesoftware bestimmt daraus die Position der Marker. Es gibt zusätzlich die Möglichkeit sogenannte \textit{Rigid bodies} zu verwenden. Dies sind Konstruktionen, die mehrere (mindestens drei) optische Marker fest miteinander verbinden. Damit ergibt sich zusätzlich zur Position auch eine Orientierung und bei der Verwendung von mehreren \textit{Rigid bodies} kann die Software diese unterscheiden. In Abbildung\,\ref{fig:RigigBody} ist ein solcher beispielhaft dargestellt.

Die \textit{Rigid bodies} können Körperteilen zugeordnet werden und so kann mit Vorwissen über die menschliche Anatomie die Bewegung rekonstruiert werden \cite{qualisys_ab_qtm-user-manual_nodate}.

\begin{figure}[h]
    \centering
    \subfloat[Kamerapositionierung \cite{noauthor_miqus_nodate}]{
        \begin{overpic}[width=0.5\textwidth,trim = 0 0 0 0]{02_grundlagen/images/Kamerapositionierung.png}
            %\put(2,25){3\,cm}
        \end{overpic}
        \label{fig:CameraPositioning}
    }
    \subfloat[Rigidbody]{
        \begin{overpic}[width=0.3\textwidth,trim = 0 0 0 0]{02_grundlagen/images/rigidbody.jpeg}
            %\put(2,25){3\,cm}
        \end{overpic}
        \label{fig:RigigBody}
    }
    \caption{Aufbau OMC: (a) zeigt eine mögliche Kameraanordnung. Der grüne Bereich ist dabei von allen Kameras einsehbar und die Marker können hier optimal lokalisiert werden. In (b) ist beispielhaft ein Rigidbody dargestellt. Dieser besteht aus einem 3D-Druckteil an dem insgesamt 4 optische Marker angebracht sind.}
    \label{fig:OMCUse}
\end{figure}
In dieser Arbeit wird zur Evaluierung ein Kamerasystem des Herstellers \textit{Qualisys AB} bestehend aus acht Miqus-Kameras \cite{noauthor_miqus_nodate} verwendet.


\subsection{Stand der Technik}
\textcolor{red}{Nicht sicher ob dieser Teil notwendig ist, vom Gefühl her eher nein.}
\label{subsec:Stateoftheart}
\subsubsection{Kommerzielle Angebote}
Im Bereich der Sensorhandschuhe gab es bis Anfang der 2020er Jahre unterschiedliche Ansätze. Dabei wurde auf verschiedene Sensoriken gesetzt, wie zum Beispiel Stretchsensoren, IMUs oder inkrementelle Messysteme. Eine Zusammenstellung verschiedener Ansätze kann hier gefunden werden\,\cite{caeiro-rodriguez_systematic_2021}. 

Optische Systeme haben sich bisher nicht durchgesetzt, da zwischen Kamera und Marker zu jedem Zeitpunkt Sichtkontakt vorhanden sein muss. Aufgrund der Anatomie der Hand können sich Finger unter Umständen gegenseitig verdecken und die Anzahl der benötigten Kameras wird zu groß und damit auch der finanzielle Aufwand.

In den letzen Jahren sind verschiedene magnetische Ansätze im kommerziellen Bereich erschienen. Ein rein magnetischer wird von der Firma \textit{Manus Meta} verfolgt, die in den neusten Versionen einen Aktuator auf dem Handrücken platzieren und jeweils Magnetfeldsensoren in den Fingerspitzen \cite{noauthor_motion_nodate},\cite{saggio_quasi-static_2025}. Laut Hersteller wird eine Endpunktlokalisierung verwendet. Eine zuvor durchgeführte Kalibrierung der Antomie ermöglicht es, aus den Positionen der Fingerspitzen Gelenkwinkel zu bestimmen.

\begin{figure}[h]
    \centering 

    \begin{overpic}[width=0.25\textwidth,trim = 0 0 0 0]{02_grundlagen/images/Manus_Glove.png}
    \end{overpic}
    \label{fig:ManusGlove}
    \caption{Kommerzieller Sensorhandschuhe: In der Abbildung ist der Sensorhandschuh von \textit{ManusMeta} zu sehen.\cite{noauthor_motion_nodate}}

\end{figure}
Das kommerzielle System verwendet jeweils ein Sensor pro Finger. Durch Kalibrierung der Handanatomie, kann aus der Sensorposition direkt die Haltung des Fingers bestimmt werden.

\subsubsection{Forschung}
In der Forschung gibt es verschiedene Arbeiten, die sich mit der Rekonstruktion von Handbewegungen basierend auf magnetischen Sensoren beschäftigen. Zum einen gibt es den Ansatz\,\cite{santoni_magik_2021}. In diesem Ansatz werden Sendespulen auf dem Handrücken und den Fingerspitzen angebracht. In einem Holzaufbau sind 24 Sensoren angebracht. Die Sendespulen werden mit einem Optimierungsalgorithmus lokalisiert. Die Hand wird als kinematische Kette beschrieben und mit einer inversen Kinematik werden die Positionen zu Winkeln übersetzt.

In einem anderen Ansatz\,\cite{ma_magnetic_2011} werden permanent Magneten an den Fingerspitzen befestigt. Um das Handgelenk ist ein Armband mit sechs 2D-Sensoren angebracht. Dieser Ansatz ist zwei geteilt, zunächst werden die Magneten auf den Fingerspitzen lokalisiert. Dafür werden die Magneten als magnetische Dipole beschrieben und mit einem Optimierungsalgorithmus eine Kostenfunktion minimiert und so die Positionen bestimmt. In einem zweiten Schritt werden mit einem iterativen Verfahren aus den Positionen Gelenkwinkel geschätzt. Dabei wird die Abbduktion/Adduktion der Finger vernachlässigt. 

\subsubsection{Zusammenfassung}
Die entwickelten magnetischen Ansätze für das Handbewegungstracking nutzen jeweils ähnliche Ansätze. In allen Ansätzen wird zunächst die Fingerspitze lokalisiert. Mit einer zuvor durchgeführten Kalibrierung wird eine inverse Kinematik angewendet und die Gelenkswinkel bestimmt. In allen Ansätzen wird das Vorwissen über die Anatomie nach der Lokalisierung in die Schätzung eingearbeitet. Die Lokalisierung kann unter Umständen sehr aufwändig sein. In dieser Arbeit soll daher ein Lokalisierungsverfahren entwickelt werden, in dem bereits während der Lokalisierung die Anatomie berücksichtigt wird.

\section{Überblick des entwickelten System}
\label{sec:SystemOverview}
In dieser Arbeit wird die Signalverarbeitung für ein Bewegungsanalysesysteme für Handbewegungen basierend auf magnetischen Feldern vorgestellt. Das System wird dafür in zwei Subsysteme unterteilt, eines für die Bewegungerkennung der Finger relativ zur Hand und eines für die Lokalisierung der Hand im Raum.

\subsection{Verwendete Sensorik}
\label{subsec:Fluxgate}
In dieser Arbeit werden Fluxgatemagnetometer(auch Förster- oder Magnetkernsonde genannt) der Firma \textit{Stefan Mayer Instruments}\,\cite{noauthor_magnetfeldsensor_nodate} verwendet, ein Bild ist in Abbildung\,\ref{fig:fluxgate} gezeigt. Diese können außerhalb einer Schirmung im Erdmagnetfeld verwendet werden und die Bandbreite von 1\,kHz ist ausreichend groß.

Ein Amplituden- und Phasenverlauf bis 3\,kHz wurde mit Stützstellen im Abstand von 100\,Hz bzw. im Abstand von 10\,Hz zwischen 500 und 1100\,Hz aufgenommen. Dieser Bereich ist höher aufgelöst, dies ist der spätere Arbeitsbereich. Der Amplituden- und Phasenverlauf ist in Abbildung \ref{fig:FluxFreq} dargestellt.
    \begin{figure}[h]
        \centering
        \scalefont{1}
            \inputTikZ{1}{02_grundlagen/images/fluxgate_Frequency.tikz}
            \label{fig:Frequenzverlauf}
        \caption{Verlauf von Sensitivität und Phase in Abhängigkeit der Frequenz}
        \label{fig:FluxFreq}

    \end{figure}
Ein Fluxgatemagnetometer besteht aus einem Ferritkern, um den zwei Spule gewickelt sind. Eine der Spulen wird mit einer Sinusansteuerung symetrisch so angesteuert, dass die Magnetisierung des Ferritkerns in beiden Richtungen die Sättigung des Ferritkerns erreicht wird. Das äußere zu messende Magnetfeld verschiebt den Mittelpunkt der Magnetisierung und in einer Richtung wird das Material stärker in die Sättigung getrieben. Mit der zweiten Spule wird die induzierte Spannung gemessen. Die Verzerrung durch die Sättigung ruft höhere harmonische der Anregungsfrequenz hervor. Die zweite Harmonische ist proportional zum äußeren Magnetfeld und wird als Maß für das Magnetfeld verwendet \cite{michalowsky_magnettechnik_2006}. In Abbildung \ref{fig:NoiseFluxgate} ist eine Rauschmessung dargestellt. Darin ist zum einen das Ansteuersignal bei ungefähr 14\,kHz zu sehen und außerdem im tieffrequenten Bereich die 50\,Hz Versorgungsnetzspannung, sowie harmonische dieser.


    
    \begin{figure}[h]
        \centering
        \subfloat[Rauschspektrum]{
            \scalefont{1} 
            \inputTikZ{1}{02_grundlagen/images/FluxGate_Noise.tikz}
            \label{fig:NoiseFluxgate} 
        } 
        \subfloat[Fluxgatemagnetometer FLC100 \cite{noauthor_magnetfeldsensor_nodate}]{
            \begin{overpic}[width=0.5\textwidth,trim = 0 0 0 0]{02_grundlagen/images/FLC_100.jpg}
            \end{overpic}
            \label{fig:fluxgate}
        }
        \caption{FLC100 mit zugehörigem Rauschspektrum: Auf der rechten Seite ist ein Foto eines FLC100 dargestellt. Dazu gehört zum einen, die Spule (Sensorelement) und die zugehörige Elektronik. Auf der linken Seite ist das Rauschspektrum zu sehen.}
        \label{fig:SensorNoise}

    \end{figure}
Für einen Prototyp wurden andere Sensoren verewndet, diese sind kleiner und können auch auf kleineren kinematischen Elementen befestigt werden. Da diese nicht zur Auswertung und Evaluieurng verwendet werden, wird hier auf eine nähere Beschreibung verzichtet.

    


\subsection{Sensor- und Aktuatoranordnung}
\label{sec:Sensoranordnung}
Bei der Sensoranordnung wird unterschieden zwischen der Lokalisierung, für die sie verwendet werden. Ein Teil der Sensoren wird zum Bewegunstracking von kinematischen Ketten verwendet. Diese werden jeweils an einem kinematischen Element befestigt. Auf dem ersten kinematischen Element ist ein magnetischer Aktuator befestigt. Um die Bewegung der kinematischen Kette aufzunehmen, wird die relative Position der Sensoren zum Aktuator bestimmt. Auf einen Handschuh angewendet, bedeutet dies, dass auf dem Handrücken magnetische Aktuatoren angebracht sind. Auf den folgenden Fingerelementen sind magnetische Sensoren angebracht.

    \begin{figure}[h]
        \centering
        \begin{overpic}[width=\textwidth,trim = 0 0 0 0]{02_grundlagen/images/Übersicht.png}
            \put(20,11){\parbox{.5in}{\centering AD-Wandler}}
            \put(47,11){\parbox{1.0in}{\centering Digitale Signalverarbeitung}}
        \end{overpic}
        \label{fig:Overview} 
        \caption{Übersicht des entwickelten Systems: In der Abbildung ist das System mit der Sensoranordnung skizziert. In der Mitte der Abbildung ist eine Hand zu sehen, auf dem Handgelenk ist eine 3D-Spule in angebracht. An der Hand sind an jedem kinematischen Element Magnetfeldsensoren (orange) angebracht. Um das Messvolumen herum ist ein Gerüst an dem insgesamt sechs Sensoren befestigt sind, jeweils ein Sensorpaar pro Dimension $x,y,z$.  Sowohl die Stromsignale (rot) und die Magnetfeldsignale (blau) werden digitalisiert und verarbeitet.}
        \label{fig:Sensoranordnung}
    \end{figure}

Die zweite Problemstellung ist die Bestimmung der absoluten Position und der Orientierung der kinematischen Kette dafür werden im Raum Sensoren angebracht. Bei diesen Sensoren werden Sensorpaare gebildet und diese werden so angebracht, dass die sensitiven Sensorachsen auf einer Geraden liegen. Die Anordnung der Sensoren ermöglicht es, die Symetrien  magnetischer Felder auszunutzen und so eine einfache Lokalisierung zu implementieren. Die Anordnung ist in Abbildung \ref{fig:Sensoranordnung} zu sehen.


      
      